\noindent\textit{2008 manuscript.}

\section{Entropy and Temperature}

\subsection*{Thermal Equilibrium}

EY : 20150821 Based on considering the physical setup of two systems that can only exchange energy between each other, that are in thermal contact, this is a derivation of temperature.

$U = U_1 + U_2$ is constant total energy of 2 systems $1,2$ in thermal contact \\
multiplicity $g(N,U)$ of combined system is 
\[
g(N,U) = \sum_{U_1 \leq U} g_1(N_1,U_1)g_2(N_2,U-U_1)
\]
The ``differential'' of $g(N,U)$ is 
\[
dg = \left( \frac{\partial g_1}{ \partial U_1 } \right)_{N_1} g_2 dU + g_1\left( \frac{\partial g_2 }{ \partial U_2 } \right)_{N_2} dU_2 = 0 
\]
EY : 20150821 This step can be made mathematically sensible by considering the exterior derivative $d$ of $g \in C^{\infty}(\Sigma)$, where $\Sigma$ is the manifold of states of the system, with local coordinates $N,U$, where $U$ happens to be a global coordinate. Then, consider a curve in $\Sigma$ s.t. it has no component in $\frac{\partial}{ \partial N}$, $\frac{\partial}{ \partial N_1}$, and this curve is a ``null curve'' so that the vector field $X\in \mathfrak{X}(\Sigma)$ generated by this curve is s.t. $dg(X)=0$.  

With $-dU_1 = dU_2$,
\[
\frac{1}{g_1} \left( \frac{\partial g_1}{ \partial U_1}\right)_{N_1} = \frac{1}{g_2} \left( \frac{\partial g_2}{ \partial U_2} \right)_{N_2} \Longrightarrow \left( \frac{ \partial \ln{g_1} }{ \partial U_1} \right)_{N_1} = \left( \frac{ \partial \ln{g_2}}{ \partial U_2} \right)_{N_2}
\]
Define
\[
\sigma(N,U) := \ln{ g(N,U)}
\]
Then
\[
\Longrightarrow \left( \frac{ \partial \sigma_1}{ \partial U_1} \right)_{N_1} = \left( \frac{ \partial \sigma_2}{ \partial U_2} \right)_{N_2}
\]



\subsection*{Temperature}

$T_1=T_2$ - temperatures of 2 systems in thermal equilibrium are equal.  

$T$ ``must be a function of $ \left( \frac{ \partial \sigma}{ \partial U} \right)_N$ \cite{CKittelHKroemer1980}.  
\[
\Longrightarrow \frac{1}{T} = k_B \left( \frac{ \partial \sigma }{ \partial U} \right)_N
\]
Experimentally, $k_B = 1.381 \times 10^{-23} \, J/K = 1.381\times 10^{-16} \, \text{ ergs}/ K$.  

Now 
\[
\begin{gathered}
\frac{1}{\tau} = \left( \frac{ \partial \sigma }{ \partial U} \right)_N
\tau = k_B T
\end{gathered}
\]


\subsection*{Problems}

\solutionhead{1} \textbf{Entropy and temperature}.  
\begin{itemize}
\item[(a)] Recall that $\frac{1}{\tau} \equiv \left( \frac{ \partial \sigma }{ \partial U} \right)_{N,V}$ and $\sigma(N,U) \equiv \log{ g(N,U)}$.  Given $g(U) = CU^{3N/2}$, 
\[
\sigma(N,U) = \log{ CU^{3N/2}} = \log{C} + \frac{3N}{2} \log{U}
\]
\[
\frac{ \partial \sigma}{ \partial U} = \frac{ 3N }{2} \frac{1}{U} = \frac{1}{ \tau} \Longrightarrow \boxed{ U = \frac{ 3N}{2} \tau  }
\]
\item[(b)] $\left( \frac{ \partial^2 \sigma}{ \partial U^2 } \right)_N < 0$ ?
\[
\left( \frac{ \partial^2 \sigma}{ \partial U^2} \right)_N = - \frac{3 N}{2} \left( \frac{1}{U^2} \right) < 0 
\]
\end{itemize}

\solutionhead{2} \textbf{Paramagnetism}.    

\[
U(s) = U_1(s_1) + U_2(s_2) = -2m B(s_1 + s_2) = -2mBs \text{ or } s = \frac{ U }{ -2 mB }
\]
i.e. potential energy $U(s) = -2s\cdot mB$.  

For $|s| \ll N$, then 
\[
g(N,s) \simeq g(N,0) \exp{ \left( -2s^2 /N \right) } = g(N,0) \exp{ \left( \frac{ - U^2 }{ 2 (mB)^2 N } \right) }
\]
\[
\sigma(N,U) = \ln{ g(N,U) } = \sigma_0 - \frac{ U^2 }{ 2m^2 B^2 } \frac{1}{N} \text{ where } \sigma_0 = \ln{ g(N,0) }
\]

\[
\frac{1}{\tau} = \left( \frac{ \partial \sigma }{ \partial U } \right)_N = \frac{ - U}{ m^2 B^2 } \frac{1}{N}
\]

What is the thermal equilibrium value of this $N$-spin system of fractional magnetization?  If $U$ denotes $\langle U \rangle$, thermal average energy, we also get the thermal average spin excess.  
\[
\langle U \rangle = \langle - 2 mB s \rangle = - 2mB \langle s \rangle
\]
\[
\Longrightarrow \tau = \frac{ m^2 B^2 N }{-U} = \boxed{ \frac{ m B N }{ 2 \langle s \rangle } }
\]


\solutionhead{3}  \textbf{Quantum harmonic oscillator}.
\begin{itemize}
\item[(a)]
Result from Ch. 1: $g(N,n) = \frac{ (N+ n -1)! }{ n! (N-1)!}$.  

Let $N-1 \to N \Longrightarrow g(N+1,n) = \frac{ (N+ n )!}{ n! N! }$.  

\[
\begin{gathered}
  \sigma(N+1,n) \equiv \ln{ g(N+1,n) } = \ln{ \frac{ (N+n)! }{ n! N! } } = \ln{ (N+n)! } - \ln{ (n!)} - \ln{ (N!) } \\
  \approx (N+n) \ln{ (N+n)} - N- n - n \ln{ n} + n - N \ln{ N } + N = \boxed{ (N+n) \ln{(N+n)} - n \ln{ n} - N\ln{N } }
\end{gathered}
\]
\item[(b)] Let $U$ denote total energy $n\hbar \omega$ of oscillators.   \\
$U = n \hbar \omega$ or $ n = \frac{ U }{ \hbar \omega}$

\[
\sigma(N, U) = (N + \frac{ U}{ \omega} ) \ln{ ( N + \frac{U }{ \omega}) } - \frac{U}{ \omega} \ln{ \frac{U}{ \omega} } - N \ln{ N}
\]

At $\tau$, $\frac{1}{\tau} = \left( \frac{ \partial \sigma}{ \partial U} \right)_N$.  
\[
\begin{gathered}
  \frac{1}{\tau} = \frac{1}{\omega} \ln{ ( N + \frac{U}{\omega})} - \frac{1}{ \omega} \ln{ \frac{ U}{ \omega} } = \frac{1}{ \omega} \ln{ \left( \frac{ N \omega }{U} + 1 \right) } \text{ or } \boxed{ U = \frac{ N \omega}{ \exp{ ( \omega / \tau) } - 1 }  }
\end{gathered}
\]

\end{itemize}

\solutionhead{4}  \textbf{The meaning of ``never.''}

Suppose $10^{10}$ monkeys.  
\begin{itemize}
\item[(a)] Hamlet represents one specific ordering of $10^{15}$ with 44 possibilities for each character.  The probability of hitting upon Hamlet from a given, random sequence is $\left( \frac{1}{44} \right)^{100000} = \boxed{ \frac{1}{44^{100000}} }$.  Given that $\log_{10}{44} = 1.64345$, then $10^{1.64345} = 44$ or $10^{-1.64345} = 44^{-1}$ so then 
\[
\left( \frac{1}{44} \right)^{100000} = 10^{-164345}
\]
\item[(b)] 
\[
(\text{age of universe})\cdot \left( \frac{ 10 \text{ keys} }{ \text{ second} } \right) = 10^{18} \, s \left( \frac{ 10 \text{ keys }}{ \text{ second } } \right) = 10^{19} \text{ keys }
\]
\[
10^{19} \text{ keys} \cdot 10^{10} \text{ monkeys} = 10^{29} \text{ keys typed out } \left( \frac{ 1 \text{ hamlet } }{ 10^5 \text{ characters } } \right) = 10^{24} \text{ possible ``Hamlets'' }
\]
From part (a), the probability that a given, random sequence is Hamlet, $10^{-164345}$
\[
(10^{29} \text{ characters})( 10^{-164345}) =  10^{-164316} 
\]

Note, I think that the probability should be $(10^{29} \text{ characters})\left( \frac{ 1 \text{ Hamlet}}{ 10^5 \text{ characters} } \right)(10^{-164345} ) = 10^{-164321}$

Since we are considering the number of ``Hamlet'', $10^5$ character sequences.  
\end{itemize}


